% Generated draft: requires mathematical and source review.
\documentclass[11pt]{article}
\usepackage[margin=25mm]{geometry}
\usepackage{fontspec}
\setmainfont{Latin Modern Roman}
\usepackage{amsmath,amssymb,mathtools,unicode-math}
\setmathfont{Latin Modern Math}
\usepackage{enumitem,needspace,xurl,hyperref}
\setlength{\parindent}{0pt}
\setlength{\parskip}{5pt}
\setlength{\emergencystretch}{3em}
\newcommand{\N}{\mathbb{N}}
\newcommand{\Z}{\mathbb{Z}}
\newcommand{\Q}{\mathbb{Q}}
\newcommand{\R}{\mathbb{R}}
\newcommand{\C}{\mathbb{C}}
\newcommand{\F}{\mathbb{F}}
\newcommand{\A}{\mathbb{A}}
\newcommand{\PP}{\mathbb{P}}
\newcommand{\E}{\mathbb{E}}
\newcommand{\eps}{\varepsilon}
\newcommand{\dd}{\,\mathrm{d}}
\begin{document}
\section*{UnsolvedMath ID: 30004854}
\textbf{Source problem code:} catalogue entry
\subsection{Definitions and mathematical statement}
The supplied record poses three related research questions about a singularity denoted by \(A\). It does not specify the class of singularities, the contact-homology theory, or the spectral sequence referred to in the original passage. The following terminology is conditional and does not supply these missing hypotheses.

For a chosen complex analytic representative \(A\), let \(A_{\mathrm{sing}}\) denote its singular locus and let \(\Delta=\{z\in\C:|z|<1\}\). A working short-arc condition on a holomorphic map \(\gamma:\Delta\to A\) is
\[
\gamma^{-1}(A_{\mathrm{sing}})=\{0\}.
\]
Write \(\operatorname{ShArc}(A)\) for the resulting space of maps, subject to the intended convention for holomorphic extension at the boundary and the topology on this space. Those conventions are not specified in the supplied passage; in particular, replacing disc maps by germs must not be assumed to preserve the requested information.

If a representative \(A\subset\C^N\) centered at the singular point has a smooth transverse link for a sufficiently small radius \(\varepsilon>0\), write
\[
L_A=A\cap S^{2N-1}_{\varepsilon},\qquad
\xi_A=TL_A\cap iTL_A.
\]
Here the intersection is required to lie in the smooth locus of \(A\), and \(\xi_A\) denotes the induced contact distribution when this construction applies. For a singularity of complex dimension \(2\), such a link is a real three-dimensional contact manifold. The record does not state hypotheses ensuring that this link construction applies.

With these unresolved choices made explicit, the supplied questions are:
\begin{enumerate}
\item Can the contact-homology algebra associated with the singularity or its contact link be computed from \(\operatorname{ShArc}(A)\)? The original phrase ``such an algebra'' has no antecedent in the supplied excerpt; its identification as a contact-homology algebra comes from the supplied cleaned statement.
\item Is there an arc-space version of the spectral sequence referred to in the source passage? The supplied excerpt does not identify that spectral sequence, its filtration, its pages, or its abutment, so no specific spectral-sequence assertion can be formulated from this record alone.
\item When \(A\) has complex dimension \(2\) and its link is a suitable contact three-manifold, can the embedded contact homology of \((L_A,\xi_A)\) be described in terms of arcs on \(A\)? The relevant arc space, coefficient choices, homology-class summands, gradings, and intended meaning of such a description are not specified.
\end{enumerate}
These are requests for constructions or computations, not a supplied conjectural isomorphism. The record does not specify whether an arc-space description should depend only on a topology, on an analytic structure, or on additional geometric data.
\subsection{Short English statement}
The problem asks whether short arcs on a singularity can compute its contact-homology algebra, whether a spectral sequence mentioned in the original report has an arc-space counterpart, and whether arcs describe embedded contact homology of the link in complex dimension \(2\). The singularity hypotheses and the referenced spectral sequence are missing, so the formulation remains conditional and requires human mathematical review.
\subsection{Sources}
\begin{itemize}
\item Local UnsolvedMath registry, supplied record with permanent ID \(30004854\), title ``Arc-Space Models of Contact Homology'', exact source code \texttt{catalogue entry}, in the Oberwolfach Reports Open Problems collection. No individual proposer is supplied; the proposed year is recorded as \(2021\).
\item Supplied report citation: \emph{Singularities (2022)}, Oberwolfach Reports, pp.\ 2573--2639, DOI \texttt{10.4171/owr/2021/47}. \url{https://doi.org/10.4171/owr/2021/47}. The supplied citation uses \(2022\), while its DOI and the recorded proposed year contain \(2021\); these details are retained without reconciliation. No pinpoint page for the question is supplied.
\item Supplied literature source: Koll\'ar--N\'emethi, ``holomorphic and short arcs'', \url{https://doi.org/10.1007/s00222-014-0530-3}. The original excerpt refers to \texttt{[KN14]}; its complete bibliography entry is not supplied.
\item Supplied literature source: ``OWR report posing the contact-homology question'', \url{https://ems.press/content/serial-article-files/46924}.
\item The original excerpt attributes embedded contact homology to \texttt{[Hut14]}. No full bibliographic entry or link for this label is supplied.
\item Literature assessment dated exactly \texttt{2026-08-22}: the supplied curation labels the problem open and states that short-arc spaces are understood for surface singularities and isolated quotient singularities, while the proposed algebraic computation of contact homology and an arc-space spectral sequence for embedded contact homology remain open research directions. This is supplied dated curation, not an independently verified assessment; the background explicitly requests independent verification of time-sensitive status and cited claims.
\item The registry labels the statement \texttt{corrected\_verified} and records a manual semantic, typographic, and LaTeX review claiming that the cleaned statement is self-contained. This is supplied review metadata, not independent mathematical verification; the missing antecedents and hypotheses identified here remain unresolved.
\end{itemize}
\end{document}
